\addtocontents{toc}{\protect\Needspace{3\baselineskip}}
\originalchapter{第三次测验复习：CMB 与暗物质阅读题}
这两组公开复习题取自 8.286 Fall 2013 Quiz 3 Review Problems。多项选择的选项也保留为文字，因其错误干扰项正好暴露物理对象混淆。历史任务用当时课程数据答题；选择题中的暗物质候选名单属于当时课堂分类，具体模型约束会随研究变化。

\src{8.286 Quiz 3 Review Problem 1(a)--(e), 原题为2007年 Quiz 3 Q1；官方复习解答。}
\begin{exercise}
(a) 判断 CMB 陈述中何者错误：扣除偶极后天空平均 $T=2.725\,\mathrm K$；扣除后 RMS 温度涨落 $\delta T/T=1.1\times10^{-3}$；大偶极主要来自观测者运动的 Doppler 效应；Penzias、Wilson 当年测得近各向同性约 $3.5\,\mathrm K$ 的额外天线温度，邻近团队解释其热宇宙遗迹意义。
(b) 大气水汽吸收短于约 $3\,\mathrm{cm}$ 的微波时，哪种历史观测方案\emph{不能凭所述地理条件}克服：高空球载 MAXIMA、南极地面 DASI、北极地面 BOOMERANG、太空 COBE/WMAP/Planck？
(c) 对 Sachs--Wolfe 与 CMB 起伏，哪项错误：光子爬出势阱变红、落下变蓝；末次散射时非重子暗物质对引力势重要；大角起伏与原初暗物质势扰动有关；角功率谱的峰仅由非重子暗物质密度起伏构成，重子本身不会产生峰？
(d) 在大质量中微子、轴子、顶夸克物质、WIMP、原初黑洞中，哪一项不是可持续存在的暗物质候选？
(e) 给两种寻找暗物质引力踪迹的观测办法。
\end{exercise}
\begin{solution}
(a) 错在 $1.1\times10^{-3}$，本题正确量级是 $1.1\times10^{-5}$；偶极运动学幅度与扣除后的原初起伏不能混同。(b) 错误方案是“北极地面 BOOMERANG”：北极海平面不能像高海拔南极干燥地点那样减小大气柱，BOOMERANG 实为从南极洲升空的球载实验；太空与高空球载另有明确避开大气的机制。(c) 错在声称峰仅由非重子暗物质造成：声学峰的振荡结构来自重子--光子流体，暗物质通过引力势与动力学参与，既不能删去重子也不能把所有大角起伏归为普通 Doppler 项。(d) 顶夸克极不稳定，不能作为长期冷暗物质粒子；其他列项是不同假设下的历史候选，是否满足全部约束需另判。(e) 例如用盘星旋转曲线与引力透镜两种独立方法；也可用星系团 virial 动力学和 CMB 功率谱模型。曲线从外盘速度推质量，透镜从光路弯曲推投影质量，误差来源不同。
\end{solution}

\src{8.286 Quiz 3 Review Problem 2(a)--(d), 原题为2009年 Quiz 3 Q1；官方复习解答。}
\begin{exercise}
(a) 给出课程口径今天 CMB 的平均温度、本星系群相对 CMB 的速度占 $c$ 的比例、扣偶极后的本征 $\delta T/T$、重子光子数比 $\eta$、末次散射面一 Hubble 长度的张角，数量级即可。
(b) 氢半电离时 $\kb T_\gamma/(13.6\,\mathrm{eV})$ 在 $1/5,1/50,10^{-3},10^{-4},10^{-5}$ 中选哪项？
(c) Sachs--Wolfe 选择题：朝我们的流体 Doppler 变红或变蓝，过密末次散射区光子爬出势阱变红或变蓝，或星系际吸收使它变红或变蓝，选对应项。
(d) 判断五句真假：暗物质同时通过引力、弱力和电磁力作用；星系团 virial 定理可估总质量；中微子被认为贡献暗物质密度的重要部分；磁单极子被认为贡献其重要部分；透镜排除 MACHO 组成全部晕，但旧课程允许最多约 $20\%$。
\end{exercise}
\begin{solution}
(a) 当年题解用 $T_0=2.725\,\mathrm K$，本星系群速度约 $630\,\mathrm{km\,s^{-1}}\simeq2\times10^{-3}c$，本征涨落约 $1.1\times10^{-5}$，$\eta\simeq6.2\times10^{-10}$，末次散射面一 Hubble 长度约 $1^\circ$。本星系群速度需扣除航天器、地球和太阳等局部速度；$\eta$ 是数密度之比，不是质量密度参数。
(b) 约 $3000\,\mathrm K$ 时 $\kb T\sim0.26\,\mathrm{eV}$，与13.6 eV之比约 $1/50$；氢半电离远晚于 $\kb T$ 直接等于结合能，因为光子多且分布有高能尾。(c) 选择“过密区出发的光子爬出引力势阱而变红”；朝我们运动的流体若仅考虑 Doppler 会蓝移，星系际吸收不是该重力项。
(d) 按原题依次假、真、假、假、真。把暗物质称“暗”不表示它\emph{不能}有任何非引力弱耦合，但题干断言普遍有电磁相互作用与通常非重子暗物质定义不符。最后一句的 $20\%$ 是旧题给定的历史概述，不能当作不带模型与天体质量范围的现时普遍上限；透镜约束依 MACHO 质量、晕模型和样本而变。
\end{solution}
